% Generated draft: requires mathematical and source review.
\documentclass[11pt]{article}
\usepackage[margin=25mm]{geometry}
\usepackage{fontspec}
\setmainfont{Latin Modern Roman}
\usepackage{amsmath,amssymb,mathtools,unicode-math}
\setmathfont{Latin Modern Math}
\usepackage{enumitem,needspace,xurl,hyperref}
\setlength{\parindent}{0pt}
\setlength{\parskip}{5pt}
\setlength{\emergencystretch}{3em}
\newcommand{\N}{\mathbb{N}}
\newcommand{\Z}{\mathbb{Z}}
\newcommand{\Q}{\mathbb{Q}}
\newcommand{\R}{\mathbb{R}}
\newcommand{\C}{\mathbb{C}}
\newcommand{\F}{\mathbb{F}}
\newcommand{\A}{\mathbb{A}}
\newcommand{\PP}{\mathbb{P}}
\newcommand{\E}{\mathbb{E}}
\newcommand{\eps}{\varepsilon}
\newcommand{\dd}{\,\mathrm{d}}
\begin{document}
\section*{UnsolvedMath ID: 20002790}
\textbf{Source problem code:} catalogue entry
\subsection{Definitions and mathematical statement}
The source poses five related research tasks: (1) isolate on the geometric trace-formula side continuous ``special'' representations, where special means larger poles of relative $L$-functions, with possible Rankin--Selberg/doubling connections attributed to Piatetski-Shapiro--Rallis and mentioning $GL_n$ or $GL_2$; (2) list precise local and global beyond-endoscopy statements, including matching; (3) describe and name Arthur's $r$-beyond-endoscopy groups and explain geometric matching; (4) explain how functional equations and an $r$-Fourier transform give an $r$-trace formula; (5) isolate hyper-Kloosterman contributions in Kuznetsov and interpret them spectrally. No field, reductive group, $L$-group representation $r$, relative datum, test space, trace normalization or exact integral identity is specified. The aid supplies two limited templates, not a complete answer to these five tasks.

The conditional analytic template is a residue/Fubini criterion. Let $(\Pi,\mu)$ be a spectral measure space, let $a(\pi)$ be a test weight, and suppose on a common punctured disc $D$ around $s=1$ a geometric meromorphic family agrees with
\[
J(s)=\int_\Pi a(\pi)\left(-\frac{L'}{L}(s,\pi,r)\right)\,\dd\mu(\pi).
\]
Assume the geometric family has no other pole in $D$, almost every $L(s,\pi,r)$ has no zero or pole in $D\setminus\{1\}$, and a contour-uniform integrable bound permits interchange of the spectral integral with integration around $\partial D$. Then the supplied criterion gives
\[
\operatorname{Res}_{s=1}J(s)
=\int_\Pi a(\pi)\bigl[-\operatorname{ord}_{s=1}L(s,\pi,r)\bigr]\,\dd\mu(\pi).
\]
A pole contributes its positive multiplicity and a zero contributes negative multiplicity. A separate radius for each spectral parameter is not the common-disc hypothesis. The statement assumes the trace identity, continuation, exclusion of other singularities and domination; it does not establish any of them for a specified automorphic problem.

The finite-field model is unconditional within its definitions. Let $\F_q$ be any finite field, $\psi$ a nontrivial additive character, $G=\F_q^\times$, and $n\geq1$. Put
\[
\operatorname{Kl}_n(a)=\sum_{x_1\cdots x_n=a}\psi(x_1+\cdots+x_n),\qquad
(H_nf)(x)=q^{-n/2}\sum_{a\in G}\operatorname{Kl}_n(a)f(x/a).
\]
Use the normalized counting inner product on functions on $G$. For a multiplicative character $\chi$, with $\tau(\chi^{-1},\psi)=\sum_{t\in G}\chi^{-1}(t)\psi(t)$, the aid proves
\[
H_n\chi=q^{-n/2}\tau(\chi^{-1},\psi)^n\chi.
\]
The convolution sum is not divided by $q-1$; that convention fixes the eigenvalue. Let $P_0f=(q-1)^{-1}\sum_{x\in G}f(x)$ be the constant projection. Since the nontrivial Gauss sums have modulus $\sqrt q$ and the trivial one is $-1$,
\[
H_n^*H_n=I-(1-q^{-n})P_0,
\]
\[
H_n=(-1)^nq^{-n/2}P_0+H_n^\circ,
\qquad (H_n^\circ)^*H_n^\circ=I-P_0.
\]
Thus the transform is unitary on the mean-zero subspace, with an explicit exceptional constant mode. For $q=2$ that subspace is zero and the statement has its usual vacuous meaning.

This is an arithmetic model of separating a special mode and of Gauss-sum/gamma-factor multiplication. It is not an isolation theorem for the global Kuznetsov formula, a construction of general $r$-Fourier transforms, a classification of Arthur's groups, or the requested complete local/global matching list. The original group and doubling-method attribution is retained as source context without inventing a more precise relative integral.
\subsection{Short English statement}
The source combines five beyond-endoscopy and Kuznetsov tasks. The aid gives a conditional way to commute a residue with a spectral integral and an exact finite-field hyper-Kloosterman transform with one quantified constant-mode defect. These templates do not supply the missing global trace identity, matching theory, groups, or automorphic interpretation.
\subsection{Sources}
\begin{itemize}
\item Local UnsolvedMath registry, permanent ID $20002790$, \texttt{catalogue entry}; \texttt{aim-representation-theory-notes.json}, \texttt{notes[61]}, source item 1 containing five numbered tasks.
\item American Institute of Mathematics, \emph{Automorphic kernel functions}, \url{https://aimath.org/pastworkshops/automorphkernelproblems.pdf}.
\item Source attributions: Piatetski-Shapiro--Rallis for the doubling-method context, Arthur for the $r$-trace formula, and the Kuznetsov/hyper-Kloosterman question. No titles, precise theorem statements, bibliography or local/global matching references are supplied.
\item Research attempt 1: partial result, problem partially solved at run, rigor label ``proved,'' novelty confidence low. Registry classification partial progress, status partially solved, statement status exact. The residue criterion and finite-field theorem are attributed to the supplied aid.
\end{itemize}
\end{document}
